% !TeX program = pdflatex
%
% SHORT paper, VERIFIED-ONLY SCOPE.
%
% RULE (River, overriding all earlier briefs): the paper reports only what the Lean
% kernel has checked.  Every theorem below carries its Lean declaration name, and every
% one was verified by the author of this file by reading
%   /home/ubuntu/fable-episode-2/zeta-math-2/lean/ZetaLucas/*.lean
% and running `lake build` (full build clean, 8671 jobs) plus `#print axioms` on each
% cited declaration (propext, Classical.choice, Quot.sound only).
%
% Nothing sourced from Mathematica sweeps, no conjectures, no fifteen-family table,
% no transfer-principle material, no (8.7) note: those live in the companion papers and
% are referred to in exactly one sentence (end of Section 1).
%
\documentclass[11pt,reqno]{amsart}

\usepackage[T1]{fontenc}
\usepackage[utf8]{inputenc}
\usepackage{amsmath,amssymb,amsthm}
\usepackage{mathtools}
\usepackage{booktabs}
\usepackage{array}
\usepackage{enumitem}
\usepackage[margin=1.15in]{geometry}
\usepackage{microtype}
\usepackage[hidelinks]{hyperref}

\theoremstyle{plain}
\newtheorem{theorem}{Theorem}[section]
\newtheorem{proposition}[theorem]{Proposition}
\newtheorem{lemma}[theorem]{Lemma}
\newtheorem{corollary}[theorem]{Corollary}

\theoremstyle{definition}
\newtheorem{definition}[theorem]{Definition}

\theoremstyle{remark}
\newtheorem{remark}[theorem]{Remark}

\newcommand{\Z}{\mathbb{Z}}
\newcommand{\Q}{\mathbb{Q}}
\newcommand{\N}{\mathbb{N}}
\newcommand{\vp}{v_p}
\newcommand{\bin}[2]{\binom{#1}{#2}}
\newcommand{\HH}[2]{H^{(#1)}_{#2}}
\newcommand{\lcm}{\operatorname{lcm}}
\newcommand{\Ln}[1]{\texttt{#1}}

\numberwithin{equation}{section}

\begin{document}

\title[Formally verified congruences for Ap\'ery-like sequences]
      {Formally verified congruences for Ap\'ery-like sequences,\\
       and an arithmetic note on the Brown--Zudilin denominators}

\author{}

\date{\today}

\subjclass[2020]{11A07, 11B65, 11Y99, 68V20}
\keywords{Ap\'ery numbers, Lucas congruences, second solution, harmonic sums,
formalisation, Lean~4, Mathlib}

\begin{abstract}
We report a Lean~4 development, checked by the Lean kernel and free of \Ln{sorry}, that
proves: (i)~the Lucas congruence $a_{ap+r}\equiv a_aa_r\pmod p$ for Ap\'ery's $\zeta(3)$
numbers, for \emph{every} prime and every $a\ge0$, together with its base-$p$ digit-product
form, and the same pair of statements for the Brown--Zudilin double sum
$Q_n=\sum_{k,l}T(n,k,l)$; (ii)~a descent lemma for truncated twisted zeta values,
$p^rK^{(r)}_\chi(y)\equiv\chi(p)K^{(r)}_\chi(\lfloor y/p\rfloor)$, and its extension to
monomials; (iii)~an abstract congruence for the \emph{second} solution of an
Ap\'ery-like recurrence, $p^wB(ap+r)\equiv\chi_pB(a)A(r)\pmod p$, under seven explicit
hypotheses on the summand and on a harmonic weight of uniform degree $w$, with the
character value $\chi_p$ carried as a parameter; (iv)~two instances of it, at $w=3$ for
Ap\'ery's $\zeta(3)$ pair and at $w=2$ for Franel's, both instantiated at $\chi_p=1$;
(v)~the closed form
$\sum_{k\le n}\bin nk^2\bin{n+k}k^2\bigl(2\HH3n-\HH3k\bigr)=b_n$, where $b_n$ is the sequence
\emph{defined} by Ap\'ery's recurrence with $b_0=0$, $b_1=6$, and, as a consequence, the
second-row congruence $p^3b_{ap+r}\equiv b_a\,a_r\pmod p$ for that classical companion
sequence itself, for every prime $p$ and all $a,r<p$; and (vi)~a small arithmetic fact taking
Brown--Zudilin's printed value $P_2=1190161/384$ as a quoted hypothesis: $d_2^5P_2\notin\Z$,
$12\,d_2^5P_2\in\Z$, and --- minimality --- any $c\in\N$ with $c\,d_2^5P_2\in\Z$ is divisible
by $12$. Each statement below is accompanied by the name of the declaration that proves it,
so that it can be checked against the repository; the paper claims nothing that the kernel
has not checked. In particular nothing about the \emph{value} $\zeta(3)$ is formalised
anywhere: this is a paper about congruences satisfied by certain sequences.
\end{abstract}

\maketitle

% ============================================================================
\section{Introduction and scope}
\label{sec:intro}

Let $a_n=\sum_{k=0}^n\bin nk^2\bin{n+k}k^2$ be Ap\'ery's numbers for $\zeta(3)$ (OEIS
\texttt{A005259}). They satisfy a congruence of Lucas type,
\begin{equation}\label{eq:arow}
   a_{ap+r}\equiv a_a\,a_r\pmod p\qquad(p\text{ prime},\ a\ge0,\ 0\le r<p),
\end{equation}
hence $a_n\equiv\prod_ia_{n_i}\pmod p$ over base-$p$ digits. This is classical
\cite{Gessel1982}, and Malik and Straub \cite{MalikStraub} prove it for all fifteen
sporadic Ap\'ery-like sequences, Straub \cite{Straub2301} sharpening it to modulus $p^2$.
Ap\'ery's irrationality proof, however, uses a \emph{second} solution $b_n$ of the same
recurrence, and for that row no congruence of Lucas type appears in the literature.

This paper reports what has been proved about the second row \emph{and machine-checked}.
The scope rule is strict and we state it once, because it is the point of the paper:

\begin{quote}
Every mathematical assertion in Sections~\ref{sec:setup}--\ref{sec:twelve} is the informal
reading of a Lean~4 declaration that the Lean kernel has checked, with no \Ln{sorry}, no
\Ln{admit}, no \Ln{native\_decide}, and no axioms beyond \Ln{propext},
\Ln{Classical.choice} and \Ln{Quot.sound}. Each is displayed with the name of its
declaration. Nothing that has been checked only numerically, or proved only on paper, is
stated here as a result.
\end{quote}

Table~\ref{tab:coverage} is the inventory. \S\ref{sec:build} says how to rebuild and
re-check the development in full.

\begin{table}[ht]
\centering\footnotesize
\caption{The development. Each module is the file
\Ln{ZetaLucas/}$\langle$module$\rangle$\Ln{.lean}, and every declaration listed lives in the
\Ln{ZetaLucas} namespace with the single exception of \Ln{Finset.sum\_range\_mul}.}
\label{tab:coverage}
\begin{tabular}{@{}llc@{}}
\toprule
module & principal declarations & \S\\
\midrule
\Ln{Core} & \Ln{choose\_digits}, \Ln{choose\_digits'}, \Ln{choose\_digits\_zero},
  \Ln{choose\_carry\_zero}, & \ref{sec:setup}\\
 & \Ln{Finset.sum\_range\_mul}, \Ln{padicValNat\_choose\_lt\_sq} & \\
\Ln{PadicBridge} & \Ln{PInt}, \Ln{PDvd}, \Ln{PCong}, \Ln{PCong.of\_zmod} & \ref{sec:setup}\\
\Ln{Apery} & \Ln{apery\_lucas}, \Ln{apery\_lucas\_digits} & \ref{sec:firstrow}\\
\Ln{BrownZudilin} & \Ln{Q\_lucas}, \Ln{Q\_lucas\_digits} & \ref{sec:firstrow}\\
\Ln{Letters} & \Ln{Letter}, \Ln{K\_descent}, \Ln{mon\_descent} & \ref{sec:LB}\\
\Ln{TheoremLB} & \Ln{theorem\_LB} & \ref{sec:LB}\\
\Ln{Instances} & \Ln{bMin}, \Ln{bMin\_lucas}, \Ln{bFranel}, \Ln{bFranel\_lucas} & \ref{sec:instances}\\
\Ln{MinimalForm} & \Ln{bApery}, \Ln{bMin\_rec}, \Ln{bMin\_eq\_bApery},
  \Ln{apery\_b\_harmonic\_closed\_form}, & \ref{sec:instances}\\
 & \Ln{bApery\_lucas} & \\
\Ln{BZFactor12} & \Ln{dlcm}, \Ln{dlcm\_two\_pow\_five}, \Ln{BZ\_P2},
  \Ln{BZ\_P2\_not\_integral}, & \ref{sec:twelve}\\
 & \Ln{BZ\_P2\_factor\_twelve}, \Ln{BZ\_twelve\_minimal} & \\
\Ln{Kummer} & seven carry-counting lemmas, headed by \Ln{padicValNat\_choose\_digits} & ---\\
\bottomrule
\end{tabular}
\end{table}

Four honest qualifications, stated here rather than buried.

\begin{enumerate}[leftmargin=1.9em,label=\textup{(\alph*)}]
\item \emph{Nothing about the value $\zeta(3)$ is formalised anywhere in this development.}
  There are no limits, no irrationality statements, and not even the convergence
  $b_n/a_n\to\zeta(3)$. Every theorem below is a congruence, an integrality statement or an
  identity between explicitly defined sequences of rationals. The word $\zeta(3)$ occurs in
  this paper only as a label for which sequences are under discussion.
\item \emph{Which sequence ``$b_n$'' means.} In Lean, $b_n$ is pinned by Ap\'ery's
  recurrence together with the initial values $b_0=0$, $b_1=6$ --- the sequence-theoretic
  characterisation --- and \emph{not} by the classical double-sum expression with its inner
  alternating sum. That the two agree is proved on paper \cite{PaperLucas} and is
  \emph{not} formalised.
\item \emph{The character is never exercised.} Both instances of the general second-row
  theorem instantiate its character parameter at $1$. The theorem carries a character, but no
  instance with a nontrivial character is proved, here or on paper.
\item \emph{One module is unused.} The \Ln{Kummer} module is a carry-counting toolkit that no
  result below consumes; it is listed for completeness.
\end{enumerate}

Finally, one sentence on what is \emph{not} here. The same programme has produced further
material --- a master depth law supported by exact sweeps over the fifteen sporadic families,
harmonic decompositions for several of them, a transfer principle relating archimedean and
Volkenborn linear forms, and denominator results for the Brown--Zudilin ladder --- none of
which is machine-checked; it is reported in the companion papers
\cite{PaperLucas,PaperSharp,PaperMap}, and nothing in this paper depends on any of it.

\subsection*{Notation} $p$ is a prime, $\vp$ the $p$-adic valuation on $\Q$ with $\vp(p)=1$,
$d_n=\lcm(1,\dots,n)$, and $\HH ry=\sum_{m=1}^ym^{-r}$. In Lean, primality is carried by the
instance \Ln{[Fact p.Prime]}, and all sums are \Ln{Finset} sums over \Ln{range}.

% ============================================================================
\section{Setup: the two bridges}
\label{sec:setup}

Two small pieces of infrastructure carry everything else: a congruence relation on $\Q$, and
a digitwise form of Lucas' theorem.

\subsection{Congruences between rationals}

Second-row statements are congruences between \emph{rational} numbers with controlled
denominators, so the ambient relation cannot be \Ln{ZMod p} equality. We use the $p$-adic
norm on $\Q$ directly.

\begin{definition}[module \Ln{PadicBridge}]\label{def:bridge}
For $x,y\in\Q$ put
\[
  \Ln{PInt}\ p\ x:\iff\lvert x\rvert_p\le1,\qquad
  \Ln{PDvd}\ p\ x:\iff\lvert x\rvert_p<1,\qquad
  \Ln{PCong}\ p\ x\ y:\iff\Ln{PDvd}\ p\ (x-y),
\]
where $\lvert\cdot\rvert_p$ is Mathlib's \Ln{padicNorm}. So \Ln{PInt} is $p$-integrality,
i.e.\ membership of the local ring $\Z_{(p)}$, and \Ln{PCong}~$p\;x\;y$ is $x\equiv y\pmod
p$ inside $\Z_{(p)}$.
\end{definition}

This relation is not vacuous: since $\lvert x-y\rvert_p<1$ forces $\lvert x-y\rvert_p\le1$,
a congruence implies $p$-integrality of either side as soon as the other is $p$-integral
(\Ln{PCong.pInt}). The module proves the expected closure properties --- sums, products by
$p$-integral constants, \Ln{Finset} sums --- and one bridge lemma, \Ln{PCong.of\_zmod}:
an equality of integers in \Ln{ZMod p} entails \Ln{PCong}~$p$ of their images in $\Q$. That
lemma is what lets the integer-valued first row be used inside rational-valued second-row
proofs.

\subsection{Lucas' theorem, digitwise}

\begin{lemma}[\Ln{Core.choose\_digits}]\label{lem:lucasstep}
Let $p$ be prime, $a,b\ge0$ and $r,s<p$. Then
$\bin{ap+r}{bp+s}\equiv\bin ab\bin rs\pmod p$.
\end{lemma}

The Lean proof specialises Mathlib's
\Ln{Choose.choose\_modEq\_choose\_mod\_mul\_choose\_div} and computes the four base-$p$
division and remainder terms explicitly. Two consequences are used repeatedly:
\Ln{choose\_digits\_zero}, that $\bin nk\equiv0\pmod p$ whenever the low digit of $n$ is
strictly below the low digit of $k$; and \Ln{choose\_carry\_zero}, the carry-annihilation
form --- if $p\le x$, $y<p$ and $x-p<y$ then $\bin xy\equiv0\pmod p$, because the digits of
$x$ are $(1,x-p)$ and those of $y$ are $(0,y)$.

The third ingredient is a combinatorial identity, \Ln{Finset.sum\_range\_mul}: for any
commutative monoid,
$\sum_{k<mn}f(k)=\sum_{i<m}\sum_{j<n}f(in+j)$. This is what splits a summation range into
base-$p$ blocks, and it is used in every proof in \S\ref{sec:firstrow} and \S\ref{sec:LB}.
Finally \Ln{padicValNat\_choose\_lt\_sq} records that for $n<p^2$ and $k\le n$ Kummer's
carry count collapses to a single indicator,
$\vp\bin nk=[\,p\le k\bmod p+(n-k)\bmod p\,]$.

% ============================================================================
\section{The first row}
\label{sec:firstrow}

\begin{theorem}[\Ln{Apery.apery\_lucas}, \Ln{Apery.apery\_lucas\_digits}]\label{thm:apery}
Let $A(n,k)=\bin nk^2\bin{n+k}k^2$ and $a_n=\sum_{k=0}^nA(n,k)$. For every prime $p$, every
$a\ge0$ and every $r<p$,
\[
   a_{ap+r}\;\equiv\;a_a\,a_r\pmod p,
   \qquad\text{and}\qquad
   a_n\;\equiv\;\prod_{i}a_{n_i}\pmod p
\]
for the base-$p$ digits $(n_i)$ of $n$.
\end{theorem}

\begin{proof}[The verified argument]
Four steps. (1)~Extend the summation range from $\{0\le k\le ap+r\}$ to the full block
$\{0\le k<(a+1)p\}$; this is an equality of integers, not a congruence, because
$A(n,k)=0$ for $k>n$. (2)~Split the block into digits by \Ln{Finset.sum\_range\_mul}.
(3)~Factor each summand modulo $p$ by Lemma~\ref{lem:lucasstep} applied to both binomials,
which gives $A(ap+r,bp+s)\equiv A(a,b)A(r,s)$. (4)~The resulting double sum is a product of
two independent sums, $\sum_b A(a,b)\cdot\sum_{s<p}A(r,s)$; the first factor is $a_a$ on the
nose, and the second is $a_r$ because the terms with $r<s<p$ carry a vanishing $\bin rs$.
The digit-product form is a strong induction on $n$ using \Ln{Nat.digits\_def'}, and it is
available for \emph{all} $a\ge0$, which is what makes the induction go through.
\end{proof}

Note the absence of hypotheses: no $p\ge5$, no $a<p$, no $p$-unit condition. The same
argument, with a two-index summand, gives the analogue for the double sum that carries the
Brown--Zudilin cellular approximations to $\zeta(5)$ \cite{BZ}.

\begin{theorem}[\Ln{BrownZudilin.Q\_lucas}, \Ln{BrownZudilin.Q\_lucas\_digits}]\label{thm:Q}
Let
\[
   T(n,k,l)=\bin{n+k}n\bin nk^2\bin{n+l}n\bin nl^2\bin{n+k+l}n,
   \qquad
   Q_n=\sum_{k,l=0}^{n}T(n,k,l).
\]
For every prime $p$, every $a\ge0$ and every $r<p$, one has
$Q_{ap+r}\equiv Q_aQ_r\pmod p$, and $Q_n\equiv\prod_iQ_{n_i}\pmod p$ over base-$p$ digits.
\end{theorem}

\begin{proof}[The verified argument]
As for Theorem~\ref{thm:apery}, with both ranges extended and both split; the summand
factorises as $T(ap+r,bp+s,cp+t)\equiv T(a,b,c)T(r,s,t)\pmod p$, the three wide binomials
$\bin{n+k}n$, $\bin{n+l}n$, $\bin{n+k+l}n$ being handled in their carrying regimes by
\Ln{choose\_carry\_zero} and \Ln{choose\_digits\_zero}. One extra step is needed
relative to the single-index case: after factorisation the four sums appear in the order
$b,s,c,t$, and a \Ln{Finset.sum\_comm} is required to bring the high-digit pair $(b,c)$ and
the low-digit pair $(s,t)$ together before the quadruple sum separates into a product.
\end{proof}

% ============================================================================
\section{The second row: harmonic letters and the abstract theorem}
\label{sec:LB}

\subsection{Letters}

\begin{definition}[module \Ln{Letters}]\label{def:letter}
For a completely multiplicative $\chi:\N\to\Z$ and $r,y\in\N$ put
\[
   K_\chi^{(r)}(y)\;=\;\sum_{m=0}^{y}\frac{\chi(m)}{m^{r}}\ \in\ \Q .
\]
The sum is over \Ln{range (y+1)}; the term $m=0$ is $\chi(0)/0^r=0$ in $\Q$ for $r\ge1$, so
the object is the usual truncated twisted zeta value $\sum_{m=1}^y\chi(m)m^{-r}$, and no
\Ln{Finset.Icc} reindexing is ever required. A \emph{letter} is a structure bundling $\chi$
together with a proof of complete multiplicativity, a positive degree $r$, and an argument
form $x:\N\to\N\to\N$; a \emph{monomial} is a \Ln{List Letter}, with $\Ln{monDeg}$ the sum of
the degrees and $\Ln{monChi}\,M\,p=\prod_{\ell\in M}\chi_\ell(p)$.
\end{definition}

We deliberately do not use Mathlib's \Ln{DirichletCharacter}, which would force a modulus and
a \Ln{ZMod}-valued target; complete multiplicativity as a hypothesis field is exactly what
the descent needs, and it admits $\chi\equiv1$.

\begin{lemma}[\Ln{Letters.K\_descent}]\label{lem:K}
Let $\chi$ be completely multiplicative, $r\ge1$ and $y\ge0$. Then
\[
   p^{r}\,K^{(r)}_\chi(y)\;\equiv\;\chi(p)\,K^{(r)}_\chi\bigl(\lfloor y/p\rfloor\bigr)
   \pmod p .
\]
\end{lemma}

\begin{proof}[The verified argument]
Split the sum at $p\mid m$. The $p$-divisible layer reproduces the right-hand side
\emph{exactly}, before any congruence is taken: reindexing $m=jp$ and using complete
multiplicativity,
$\sum_{j\le\lfloor y/p\rfloor}\chi(jp)(jp)^{-r}=\chi(p)p^{-r}K^{(r)}_\chi(\lfloor
y/p\rfloor)$. The remaining layer has $p$-integral denominators, so after multiplication by
$p^{r}$ with $r\ge1$ it is $p$-divisible. Hence the difference of the two sides is
$p^r$ times a $p$-integral number.
\end{proof}

Lemma~\ref{lem:K} is where the character enters, and it is the only place. Multiplying the
lemma over the letters of a monomial gives \Ln{mon\_descent}: if every argument satisfies
$\lfloor x(n,k)/p\rfloor<p$, then
\[
   p^{\Ln{monDeg}M}\cdot M(n,k)\;\equiv\;\Ln{monChi}\,M\,p\cdot M^{\downarrow}(n,k)\pmod p,
\]
where $M^\downarrow$ is $M$ with every argument $x$ replaced by $\lfloor x/p\rfloor$. The
Lean proof is an induction over the list, using that each factor is $p$-integral so that the
cross terms in the product of two congruences are absorbed.

\subsection{The abstract second-row theorem}

Fix a summand $S:\N\to\N\to\Z$, a finite index set $J$, coefficients $c:J\to\Q$ and monomials
$\mathrm{mon}:J\to\Ln{List Letter}$, and put
\[
   w(n,k)=\sum_{j\in J}c_j\,\mathrm{mon}_j(n,k),\qquad
   A(n)=\sum_{k\le n}S(n,k),\qquad
   B(n)=\sum_{k\le n}S(n,k)\,w(n,k),
\]
the last two being \Ln{Arow} and \Ln{Brow} in the development. The theorem below carries a
character value $\chi_p$ as a parameter, and since both instances proved here set it to $1$
it is worth saying why it is there: in the companion work \cite{PaperLucas} the untwisted
form of this congruence is found to fail for seven of the fifteen sporadic Ap\'ery-like
families, at precisely the primes where a quadratic character takes the value $-1$, and
Lemma~\ref{lem:K} --- formalised here as \Ln{K\_descent} --- is the mechanism that produces
the compensating factor. We state that as motivation for the shape of the theorem, and claim
none of it here.

\begin{theorem}[\Ln{TheoremLB.theorem\_LB}]\label{thm:LB}
Let $p$ be prime and let $w\in\N$, $\chi_p\in\Z$. Assume
\begin{description}[leftmargin=3.2em,style=sameline]
\item[\textup{(H1)}] $S(ap+r,bp+s)\equiv S(a,b)S(r,s)\pmod p$ for \emph{all} $a,b\ge0$ and
  all $r,s<p$;
\item[\textup{(H2)}] $S(n,k)=0$ whenever $n<k$;
\item[\textup{(H3)}] for all $r,s<p$ and $b\le a$ with $S(r,s)\not\equiv0\pmod p$, every
  argument form $x$ occurring in the weight satisfies
  $\lfloor x(ap+r,bp+s)/p\rfloor=x(a,b)$;
\item[\textup{(H4)}] every argument form satisfies $x(n,k)\le n$ whenever $k\le n$;
\item[\textup{(H4c)}] every coefficient $c_j$ is $p$-integral;
\item[\textup{(Hw)}] every monomial has total degree exactly $w$;
\item[\textup{(H5)}] $\Ln{monChi}\,\mathrm{mon}_j\,p=\chi_p$ for every $j\in J$.
\end{description}
Then for all $a<p$ and $r<p$,
\[
   p^{w}\,B(ap+r)\;\equiv\;\chi_p\,B(a)\,A(r)\pmod p .
\]
\end{theorem}

\begin{proof}[The verified argument]
Write $n=ap+r$, so $n<p^2$ and $\lfloor n/p\rfloor=a$. By (H4) every argument of the weight
at $(n,k)$ is $\le n<p^2$, so its descent $\lfloor x/p\rfloor$ is $<p$; with (H4c) this makes
$p^ww(n,k)$ $p$-integral for \emph{every} $k$, and $w(a,b)$ $p$-integral as well. Extend the
range of $B(n)$ to the full block $\{0\le k<(a+1)p\}$ --- an equality, by (H2) --- and split
it into digits. On the layer where $S(r,s)\equiv0\pmod p$, (H1) makes each term $\equiv0$,
because it is $p$ times a $p$-integral number. On the complementary layer, (H3) identifies
the descended weight with $w(a,b)$, \Ln{mon\_descent} and (Hw), (H5) give
$p^ww(n,k)\equiv\chi_pw(a,b)$, and the double sum factors as
$\bigl(\sum_{b\le a}S(a,b)w(a,b)\bigr)\bigl(\sum_{s<p}S(r,s)\bigr)=B(a)A(r)$, the second
factor again by (H2).
\end{proof}

\begin{remark}[what the formal hypotheses cost, and save]\label{rem:hyps}
Two features of the hypothesis list are worth flagging for anyone instantiating it. (H1) is
an \emph{unconditional} congruence, required for all $a,b$ rather than only on a surviving
set; in the carrying regime both sides vanish modulo $p$, so the usual case dichotomy is
absorbed, and in consequence no product-region argument occurs anywhere in the formal proof:
the double sum factors because a full block splits, and the extension to the full block is an
equality of integers. (H3) is conditioned on $S(r,s)\not\equiv0\pmod p$, which is exactly
what excludes a borrow in an argument such as $n-k$: an instance \emph{derives} the
no-borrow condition instead of describing the surviving set. (H5) asks that the product of
the letters' characters at $p$ be the same for every monomial, which is weaker than
requiring a fixed number of twisted letters per monomial and permits mixed conductors.
\end{remark}

% ============================================================================
\section{Two instances, and a closed form}
\label{sec:instances}

\subsection{The weight-three instance}

\begin{definition}[\Ln{Instances.bMin}]\label{def:bmin}
$\displaystyle B_{\min}(n)=\sum_{k=0}^{n}\bin nk^2\bin{n+k}k^2\bigl(2\HH3n-\HH3k\bigr)$.
\end{definition}

\begin{theorem}[\Ln{Instances.bMin\_lucas}]\label{thm:bmin}
For every prime $p$ and all $a,r<p$,
\[
   p^3\,B_{\min}(ap+r)\;\equiv\;B_{\min}(a)\,a_r\pmod p .
\]
\end{theorem}

\begin{proof}[The verified argument]
Theorem~\ref{thm:LB} with $S(n,k)=A(n,k)$, $J$ of size two, coefficients $(2,-1)$, monomials
$[\HH3{\,\cdot\,}]$ at the arguments $x(n,k)=n$ and $x(n,k)=k$, degree $w=3$ and $\chi_p=1$.
(H1) is the digit factorisation of $A$ established for Theorem~\ref{thm:apery}; (H2) is
$\bin nk=0$ for $k>n$; (H3) holds because $\lfloor(ap+r)/p\rfloor=a$ and
$\lfloor(bp+s)/p\rfloor=b$; (H4) because both arguments are $\le n$; (H4c) because the
coefficients are the integers $2,-1$.
\end{proof}

There is no hypothesis $p\ge5$ and no hypothesis $a\ge1$ --- both of which the informal
analysis of this instance had carried. The first is removable precisely because the
coefficients of the minimal weight are integers, so (H4c) is vacuous; the second because
$B_{\min}(0)=0$ and $p^3B_{\min}(r)$ is $p$-divisible for $r<p$.

\subsection{The weight-two instance}

\begin{theorem}[\Ln{Instances.bFranel\_lucas}]\label{thm:franel}
Let $f_n=\sum_k\bin nk^3$ be the Franel numbers \textup{(OEIS \texttt{A000172})} and
\[
   B_{\mathbf F}(n)=\sum_{k=0}^{n}\bin nk^3
     \Bigl(\tfrac14\HH2k+\tfrac34H_k^2-\tfrac34H_kH_{n-k}\Bigr).
\]
For every prime $p\ne2$ and all $a,r<p$, one has
$p^2B_{\mathbf F}(ap+r)\equiv B_{\mathbf F}(a)f_r\pmod p$.
\end{theorem}

\begin{proof}[The verified argument]
Theorem~\ref{thm:LB} with $S(n,k)=\bin nk^3$, three monomials --- $\HH2k$, $H_kH_k$ and
$H_kH_{n-k}$ --- coefficients $\tfrac14,\tfrac34,-\tfrac34$, degree $w=2$ and $\chi_p=1$.
The only hypothesis needing care is (H3) for the argument $n-k$: it is here that the side
condition $S(r,s)\not\equiv0\pmod p$ does its work, giving $\bin rs\ne0$, hence $s\le r$,
hence $\lfloor(n-k)/p\rfloor=a-b$ with no borrow. The hypothesis $p\ne2$ enters only through
(H4c), the coefficient denominators being $4$.
\end{proof}

\subsection{The closed form}

Theorem~\ref{thm:bmin} is a statement about the explicit sum $B_{\min}$. The following
identifies that sum with Ap\'ery's companion sequence.

\begin{definition}[\Ln{MinimalForm.bApery}]\label{def:bapery}
$b:\N\to\Q$ is defined by $b_0=0$, $b_1=6$ and
$(n+1)^3b_{n+1}=P(n)b_n-n^3b_{n-1}$ with $P(n)=34n^3+51n^2+27n+5$; the leading coefficient
never vanishes, so this determines $b_n$ uniquely.
\end{definition}

This is the sequence-theoretic characterisation of Ap\'ery's companion: the recurrence
together with two initial values. It is \emph{not} the classical double-sum expression, in
which the weight attached to $\bin nk^2\bin{n+k}k^2$ is
$\HH3n+\sum_{m\le k}(-1)^{m-1}\bigl(2m^3\bin nm\bin{n+m}m\bigr)^{-1}$; that the two agree is
proved on paper in \cite{PaperLucas} and is not formalised. Every occurrence of $b_n$ below
means Definition~\ref{def:bapery}.

\begin{theorem}[\Ln{MinimalForm.bMin\_eq\_bApery} and
\Ln{apery\_b\_harmonic\_closed\_form}]\label{thm:closed}
For every $n\ge0$,
\[
   \sum_{k=0}^{n}\bin nk^2\bin{n+k}k^2\bigl(2\HH3n-\HH3k\bigr)\;=\;b_n .
\]
\end{theorem}

\begin{proof}[The verified argument]
Both sides satisfy the same order-two recurrence and agree at $n=0,1$; since the leading
coefficient $(n+1)^3$ never vanishes, a two-step induction concludes. That $B_{\min}$
satisfies the recurrence (\Ln{bMin\_rec}) is the substance. Everything is driven by four
multiplicative absorption identities (\Ln{l0a}--\Ln{l0d}) which, writing $n=m+1$ and
$\Phi(m,k)=\bin{m+2}k^2\bin{m+k}k^2/\bigl((m+1)^2(m+2)^2\bigr)$, express $A(m,k)$,
$A(m+1,k)$, $A(m+2,k)$ and $\Phi(m,k+1)(k+1)^4$ as $\Phi(m,k)$ times explicit quadratic
factors. These hold for \emph{all} $k\ge0$, including the degenerate range $k>n$, so no
$0/0$ and no boundary case ever arises, and after them every step is $\Phi(m,k)$ times a
polynomial identity in $\Q[m,k]$ discharged by \Ln{ring}. Two certificates are then
exhibited: a Zeilberger certificate $G=\Phi k^4T$ with
$T(n,k)=4(2n+1)(2k^2-3k-4n^2-4n)$, and a Gosper certificate
$K=\Phi k^4U/\bigl(n^2(n+1)^2\bigr)$ with $U(n,k)=-4(k-1)(2n+1)(2n^2+2n+1-k)$ which absorbs
the weight correction. They combine into a \emph{single} telescope: with $Gd=\Phi kT$ and
\[
   \Psi(m,k)=\bigl(2\HH3{m+1}-\HH3k\bigr)G(m,k)+K(m,k)+Gd(m,k),
\]
the $k$-summand of $(Lb)(m+1)$ is exactly $\Psi(m,k+1)-\Psi(m,k)$ (\Ln{star}), so
\Ln{Finset.sum\_range\_sub} collapses the whole sum to $\Psi(m,m+3)-\Psi(m,0)=0$. There is no
Abel summation and no boundary bookkeeping beyond the two evaluations
$\Psi(m,0)=\Psi(m,m+3)=0$.
\end{proof}

\begin{remark}[what Theorem~\ref{thm:closed} does and does not say]\label{rem:closedscope}
What is verified is an identity between an explicit binomial-harmonic sum and the solution of
an explicit recurrence. That $b_n/a_n\to\zeta(3)$, that $b_n$ coincides with Ap\'ery's
classical nested-sum expression, and that $d_n^3b_n\in\Z$ are \emph{not} formalised and are
not asserted here.
\end{remark}

Combining the last two theorems removes the explicit sum from the statement altogether, and
gives the congruence in the form a reader of \S\ref{sec:intro} would want it: a statement
about the solution of Ap\'ery's recurrence.

\begin{theorem}[\Ln{MinimalForm.bApery\_lucas}]\label{thm:bapery}
Let $b$ be as in Definition~\ref{def:bapery} and $a_n$ as in Theorem~\ref{thm:apery}. For
every prime $p$ and all $a,r<p$,
\[
   p^3\,b_{ap+r}\;\equiv\;b_a\,a_r\pmod p .
\]
\end{theorem}

\begin{proof}[The verified argument]
Theorem~\ref{thm:bmin} transported along Theorem~\ref{thm:closed}: rewrite both occurrences
of $b$ as $B_{\min}$ and apply \Ln{bMin\_lucas}.
\end{proof}

\begin{remark}[status in the literature]\label{rem:lit}
The recurrence of Definition~\ref{def:bapery} is Ap\'ery's, and that his companion sequence
satisfies it is classical; a complete creative-telescoping proof is given by Schneider
\cite{Schneider2007}, and the identity goes back to the reorganisation of Ap\'ery's proof
reported by van der Poorten \cite[\S8]{vdPoorten1979}. What we have not been able to locate
in the literature is the weight $2\HH3n-\HH3k$ of Theorem~\ref{thm:closed} --- in particular
it is not among the closed forms recorded for OEIS \texttt{A059415}. The nearest published
relative we found is Nesterenko's Lemma~1 \cite{Nesterenko1996}, which amounts to the longer
weight $\HH3k+(2H_k-H_{n+k}-H_{n-k})\HH2k$.
\end{remark}

% ============================================================================
\section{An arithmetic note on the Brown--Zudilin denominators}
\label{sec:twelve}

Brown and Zudilin \cite{BZ} construct cellular approximations to $\zeta(5)$ and record, as an
experimental observation, the inclusions $Q_n,\;d_n^2d_{2n}\hat P_n,\;d_n^5P_n\in\Z$ for the
totally symmetric member of their family. They also print the value $P_2=1190161/384$.

The following is an arithmetic consequence of that printed numeral and nothing else. We
emphasise what is and is not formalised: the module \Ln{BZFactor12} takes
$P_2=1190161/384$ as a \emph{quoted rational literal}, a hypothesis of the file; no cellular
integral, no period, and no part of the construction that produced the number is formalised
or verified. If the printed value were a typographical error, the statements below would say
nothing about the corrected value.

\begin{theorem}[\Ln{BZ\_P2\_not\_integral}, \Ln{BZ\_P2\_factor\_twelve},
\Ln{BZ\_twelve\_minimal}]\label{thm:twelve}
Let $d_n=\lcm(1,\dots,n)$ be computed from \Ln{Nat.lcm} and let $P_2=1190161/384$. Then
$d_2^5=32$ and
\[
   d_2^5\,P_2=\frac{1190161}{12}\notin\Z,
   \qquad
   12\,d_2^5\,P_2=1190161\in\Z,
\]
and moreover $12$ is minimal: every $c\in\N$ with $c\,d_2^5P_2\in\Z$ is divisible by $12$.
\end{theorem}

\begin{proof}[The verified argument]
$d_2=\lcm(1,2)=2$ and $d_2^5=32$ are computed by \Ln{decide} from the definition of
$d_n$ by recursion on \Ln{Nat.lcm}, not hard-coded; $32\cdot1190161/384=1190161/12$ by
\Ln{norm\_num}. If $1190161/12$ were an integer $z$ then $1190161=12z$ in $\Z$, which
\Ln{omega} refutes. For minimality, $c\cdot1190161/12=z$ gives $1190161c=12z$ in $\Z$, and
since $1190161\equiv1\pmod{12}$ one gets $12\mid c$; the Lean proof exhibits the witness
$c\cdot(-99180)+z$ directly.
\end{proof}

So with the printed $P_2$, the inclusion $d_n^5P_n\in\Z$ fails at $n=2$, and fails by exactly
$12=2^2\cdot3$: no smaller multiplier repairs it, and $12$ does. We record it in this form
because it is exactly the part of the matter that a kernel can settle. Whether $12$ is the
right global constant for all $n$, and where its two prime factors come from, is a
computational and structural question treated in \cite{PaperSharp}; nothing about it is
claimed here.

% ============================================================================
\section{Building and checking the development}
\label{sec:build}

The development is $2107$ lines across the ten \Ln{ZetaLucas/} modules ($2117$ including the
root import file), and depends only on Mathlib \cite{Mathlib,Lean4}. The toolchain is pinned
in \Ln{lean-toolchain} and the Mathlib revision in \Ln{lake-manifest.json}:

\begin{quote}\ttfamily\small
leanprover/lean4:v4.33.0-rc1\\
mathlib rev cd580e54f1a6b46063824e80cec92f64692cbe78\\[4pt]
lake exe cache get\\
lake build
\end{quote}

\noindent
With Mathlib prebuilt this replays in seconds and reports
\Ln{Build completed successfully}. Three independent checks establish the claim made in
\S\ref{sec:intro}, and a reader is invited to run all three:

\begin{enumerate}[leftmargin=1.9em,label=\textup{(\arabic*)}]
\item \emph{No holes.} \Ln{grep -rn "sorry\textbackslash|admit\textbackslash|native\_decide"
  ZetaLucas/} returns only one hit, inside a prose comment in \Ln{Kummer.lean}
  (``two digits admit at most one carry''). No declaration in the development is stated with
  \Ln{sorry}, and no proof uses \Ln{native\_decide}, so no statement rests on compiled-code
  evaluation.
\item \emph{No extra axioms.} \Ln{\#print axioms} on each declaration named in this paper
  returns \Ln{[propext, Classical.choice, Quot.sound]} and nothing else --- in particular no
  \Ln{sorryAx}.
\item \emph{The statements are the ones printed here.} Each theorem above cites its
  declaration; the Lean statements are short and were written to be read.
\end{enumerate}

Two small points of hygiene. The development contains a handful of \Ln{\#eval} normalisation
checks (Ap\'ery and Franel values, the first terms of $B_{\min}$ and $B_{\mathbf F}$, and
digitwise sweeps of the two instances at $p=5,7$). These are \emph{evaluations, not kernel
proofs}, they are marked as such in the source, and no statement in this paper depends on
them. Separately, five anonymous \Ln{example}s \emph{are} kernel-checked: in \Ln{Apery.lean},
by \Ln{decide}, that $a_0,\dots,a_5=1,5,73,1445,33001,819005$ and that
$a_{5a+r}\equiv a_aa_r\pmod5$ for all $a,r<5$; and in \Ln{MinimalForm.lean}, by
\Ln{norm\_num}, that $b_2=351/4$ and $b_3=62531/36$ from Definition~\ref{def:bapery} and that
$B_{\min}(2)=351/4$ from the sum.

\section*{Acknowledgements}

% WORDING FLAGGED -- please review this sentence before circulation.
This work was carried out in a human--AI collaboration: the statements were drafted and the
Lean development written in a loop between the authors and language-model agents, with every
theorem in this paper checked against the repository by rebuilding it and printing its
axioms.

% ============================================================================
\begin{thebibliography}{99}
\small

\bibitem{Beukers2211}
F.~Beukers, \emph{$p$-Linear schemes for sequences modulo $p^r$}, \texttt{arXiv:2211.15240}.

\bibitem{BZ}
F.~Brown and W.~Zudilin, \emph{On cellular rational approximations to $\zeta(5)$},
\texttt{arXiv:2210.03391}.

\bibitem{Gessel1982}
I.~M.~Gessel, \emph{Some congruences for Ap\'ery numbers}, J.\ Number Theory \textbf{14}
(1982), no.~3, 362--368.

\bibitem{MalikStraub}
A.~Malik and A.~Straub, \emph{Divisibility properties of sporadic Ap\'ery-like numbers},
\texttt{arXiv:1508.00297}.

\bibitem{Mathlib}
The mathlib Community, \emph{The Lean mathematical library}, in: CPP~2020, ACM, 367--381;
\texttt{arXiv:1910.09336}.

\bibitem{Lean4}
L.~de~Moura and S.~Ullrich, \emph{The Lean 4 theorem prover and programming language},
in: CADE~28, Lecture Notes in Comput.\ Sci.\ \textbf{12699}, Springer, 2021, 625--635.

\bibitem{Nesterenko1996}
Yu.~V.~Nesterenko, \emph{A few remarks on $\zeta(3)$}, Mat.\ Zametki \textbf{59} (1996),
no.~6, 865--880; English transl., Math.\ Notes \textbf{59} (1996), 625--636, Lemma~1.

\bibitem{Schneider2007}
C.~Schneider, \emph{Ap\'ery's double sum is plain sailing indeed}, Electron.\ J.\ Combin.\
\textbf{14} (2007), \#N5, 3 pp.

\bibitem{Straub2301}
A.~Straub, \emph{Gessel--Lucas congruences for sporadic sequences},
\texttt{arXiv:2301.12248}; Monatsh.\ Math.\ (2023).

\bibitem{vdPoorten1979}
A.~van der Poorten, \emph{A proof that Euler missed\dots\ Ap\'ery's proof of the
irrationality of $\zeta(3)$}, Math.\ Intelligencer \textbf{1:4} (1978/79), 195--203, \S8.

\medskip
\noindent\emph{Companion papers of the same programme. None of their contents is
machine-checked, and nothing in the present paper depends on them:}

\bibitem{PaperLucas}
\emph{$\chi$-twisted Lucas congruences for the second solutions of Ap\'ery-like recurrences}
(companion paper).

\bibitem{PaperSharp}
\emph{Sharp denominators for the Brown--Zudilin cellular approximations to $\zeta(5)$}
(companion paper).

\bibitem{PaperMap}
\emph{A transfer principle for $p$-adic hypergeometric linear forms, and a margin map}
(companion paper).

\end{thebibliography}

\end{document}
